\documentclass[mestrado]{UnB-PPGT}

% Metadados obrigatórios para teste
\titulo{Teste dos Comandos para Elementos Comuns}
\autor{João}{Silva}
\orientador{Prof. Dr. Maria Santos}{UnB}
\datadefesa{15}{12}{2024}
\coordenador{Prof. Dr. Carlos Oliveira}{UnB}
\membrobanca{Prof. Dr. Ana Costa}{UnB}
\membrobanca{Prof. Dr. Pedro Lima}{UFMG}

% Palavras-chave para teste
\palavraschave{teste, elementos, LaTeX, PPGT}
\keywords{test, elements, LaTeX, PPGT}

\begin{document}

\pretextual

% Teste dos ambientes pré-textuais
\begin{resumo}
Este documento testa os comandos implementados para elementos comuns no template PPGT.
\end{resumo}

\begin{abstract}
This document tests the implemented commands for common elements in the PPGT template.
\end{abstract}

\textual

\chapter{Teste dos Comandos de Elementos Comuns}\label{cap:teste}

\section{Teste de Figuras}

% Teste do comando básico de figura (existente)
% \figura[positioning]{image_file}{caption}{label}{includegraphics_options}
\figura{img/positivo_cor.txt}{Esta é uma figura de teste usando o comando básico}{fig:teste1}{width=0.8\textwidth}

% Teste de figura com largura personalizada (novo comando)
\figuracom{img/positivo_cor.txt}{0.5\textwidth}{Figura com largura personalizada}{fig:teste2}

% Teste de referência a figuras (usando comandos existentes)
Como podemos ver na \refFig{fig:teste1} e na \refFig{fig:teste2}, os comandos funcionam corretamente.

\section{Teste de Tabelas}

% Teste do comando básico de tabela (existente)
% \tabela{caption}{label}{column_spec}{table_content}
\tabela{Tabela de teste usando comando básico}{tab:teste1}{|c|c|c|}{
Coluna 1 & Coluna 2 & Coluna 3 \\
Dados 1 & Dados 2 & Dados 3 \\
Mais dados & Mais dados & Mais dados
}

% Teste de tabela PPGT com formatação específica (novo comando)
\tabelappgt{|l|r|c|}{
Item & Valor & Status \\
Item A & 100 & Ativo \\
Item B & 200 & Inativo \\
Item C & 150 & Ativo
}{Tabela com formatação específica PPGT}{tab:teste2}

% Teste de referência a tabelas (usando comandos existentes)
Os dados apresentados na \refTab{tab:teste1} e na \refTab{tab:teste2} demonstram o funcionamento dos comandos.

\section{Teste de Equações}

% Teste de equação básica (comando existente)
% \equacao{label}{equation_content}
\equacao{eq:einstein}{E = mc^2}

% Teste de sistema de equações (novo comando)
\sistema{
x + y &= 10 \\
2x - y &= 5
}{eq:sistema}

% Teste de equação em caixa (novo comando)
\equacaocaixa{F = ma}{eq:newton}

% Teste de referências a equações (usando comandos existentes)
A famosa \refEq{eq:einstein} de Einstein, o \refEq{eq:sistema} linear, e a \refEq{eq:newton} de Newton.

\section{Teste de Referências Cruzadas}

Este capítulo (\refCap{cap:teste}) demonstra o uso de:
\begin{itemize}
\item Figuras: \refFig{fig:teste1} e \refFig{fig:teste2}
\item Tabelas: \refTab{tab:teste1} e \refTab{tab:teste2}  
\item Equações: \refEq{eq:einstein}, \refEq{eq:sistema}, e \refEq{eq:newton}
\item Referências abreviadas: \reffigabrev{fig:teste1}, \reftababrev{tab:teste1}, \refeqabrev{eq:einstein}
\end{itemize}

\chapter{Conclusão}\label{cap:conclusao}

Os comandos para elementos comuns foram implementados com sucesso, conforme demonstrado no \refCap{cap:teste}.

\end{document}